\documentclass[11pt,a4paper]{article}
% Provenance: coherent_evidence Python reductions and lfm_coherent_* Rust tools.
% All numerical values below are exported from retained HDF5 products.
\usepackage[margin=20mm]{geometry}
\usepackage{amsmath,booktabs,graphicx,float,hyperref}
\usepackage[T1]{fontenc}
\usepackage{mathptmx}
\hypersetup{colorlinks=true,urlcolor=blue,citecolor=blue}
\setlength{\parindent}{0pt}\setlength{\parskip}{5pt}
\input{metrics.tex}
\begin{document}
\begin{center}{\Large Evidence audit: enhanced SYISR EPB echoes}\\[5pt]
1 March 2026, 11:45--13:45 UTC (19:45--21:45 Beijing)\\
Juha-Pekka Vierinen \hfill 21 November 2026\end{center}
\textbf{Result.} Beam 9 (azimuth $358^\circ$, elevation $60^\circ$) contains a
persistent enhanced return whose power is inconsistent with ordinary thermal
ISR under the calibration tested here. The peak apparent ISR density is
$\peakNe$ m$^{-3}$ at \peakHeight~km and 12:48:30 UTC. It is approximately
\iriExcess{} times the local IRI density. This is strong evidence for a
nonthermal or coherent return, consistent with field-aligned irregularities.
The calibration and scattering tests below distinguish that conclusion from
complete proof of the specific scattering mechanism or the absence of sidelobe
and instrumental contamination.

\section*{1. L2 density anchor and power-to-density calculation}
The power product is an arithmetic mean over individual pulses in 60 s bins,
after linear averaging to approximately 6 km range cells and per-pulse
far-range noise subtraction. Signed values are retained. Pulse medians are
used only to diagnose impulsive contamination; they are not silently treated
as calibrated mean power. All eleven beams use the same unit-energy filter,
432 MHz frequency, 479 $\mu$s recorded chirp, 4 MHz sampling and bandwidth,
and 12 ms transmitter IPP.

Two quiet zenith LP-L2 profiles at 11:02:45 and 11:19:46 supply the density
anchor. The calibration uses 300--450 km, \calibrationCount{} matched points,
\texttt{Fitstate}=0, reported density error at most 50\%, positive LFM signal
above 0.2 times the noise, and an impulsiveness screen based on mean versus
median. No upper density threshold is imposed. The standard ion-line relation
\cite{Yue2023} gives the conditional apparent density
\begin{equation}
 N_{e,\mathrm{app}}=C_b r^2(P_b-N_b)
             \left[1+(T_e/T_i)_{\mathrm{IRI}}\right].
\end{equation}
IRI-2020 provides \emph{assumed} temperatures at WGS84 scattering locations;
it does not provide an independent measured density. The zenith coefficient
is \coefficient{} for $r$ in km and power in the recorded voltage-squared
units, with a 10th--90th percentile
range \coefficientLow--\coefficientHigh{} across calibration points. This
spread is not a calibrated uncertainty interval. Quiet matched-beam control
coefficients span only 1 to \controlMax{} times the zenith value. For beam 9,
the nominal transfer is $C_9=C_{\rm zenith}/\sin60^\circ$, an ideal projected
aperture approximation for volume-filling thermal scatter. A separate factor
10 decrease and increase in inferred density tests gain-transfer sensitivity;
these are stress cases, not confidence bounds.

The LP-L2 file has approximately 17 minute sampling and contains isolated
very large estimates even where its numerical quality flag is successful.
They are retained in the audit, not used as evidence of enormous ambient
plasma density. The quiet F-region plateau and robust calibration prevent
those isolated values from dominating the anchor. The absolute calibration
history of this supplied L2 product has not yet been independently verified.
L2 UTC is registered to the original LP clock: its first timestamp is
28 February 23:59:51, matching 1 March 07:59:51.015350 Beijing in the first
LP record after the eight-hour conversion. Event LP file starts register
within 24 s of the L2 sequence; this does not imply simultaneous profile
integration or remove the approximately 17 minute sampling limitation.
SYISR ion-line calibration can vary with altitude, time and beam direction
\cite{Hao2022,Yue2023}.

\clearpage
\begin{figure}[H]\centering
\includegraphics[width=\linewidth]{figures/coherent_evidence.pdf}
\caption{Performed power and calibration diagnostics. (a) Calibration
coefficients from the two quiet zenith profiles: teal is 11:02:45 and magenta
11:19:46; black dashed is their pooled median. (b) Beam 9 apparent thermal-ISR
density, using the stated gain and IRI temperature assumptions. Low and
nonpositive power is shown with the lowest colour; signed estimates remain
in the data. Short bright spikes elsewhere are retained. (c) Maximum apparent
density among persistent strong cells (solid magenta), its factor 10 lower
stress case (dashed magenta), the 432 MHz plasma cutoff density (black dashed),
and IRI peak density in this beam (teal). (d) Far-range noise relative to its
own quiet median: beam 9 is magenta and zenith teal. Apparent density is a
null-hypothesis diagnostic, not an adopted physical density retrieval.}
\end{figure}

\section*{2. Why this is more than one bright pixel}
There are \strongCells{} strong persistent time--height cells in
\strongMinutes{} one-minute bins from \firstStrong{} to \lastStrong{} UTC,
over approximately 276--371 km. The descriptive selection requires mean
signal above 100 times noise and pulse median above half the mean.
At the maximum, the pulse median/mean ratio is \peakPersistence{}, while
large impulsive maxima in other beams often have ratios of $10^{-7}$--$10^{-4}$.
The feature therefore differs from an isolated hard-target spike.
The central ray is \peakAspect$^\circ$ from local $B$-perpendicular at the
maximum; finite beam width includes directions nearer perpendicular.

Beam 9 far-range noise changes by only \beamNoise{} between the quiet and
enhanced intervals (zenith: \zenithNoise). A simple receiver gain jump large
enough to explain the echo would also strongly affect receiver noise, which
is not observed. At the selected range, median correct/wrong-chirp power
changes from \codeQuiet{} in the quiet record to \codeEvent{} during the
event, supporting a localized return of the transmitted waveform rather
than a generic broadband increase. Exact raw-voltage extrema occur only once
or twice in each audited component, with no obvious digital plateau.
The available voltages are already beamformed: ADC rails and per-element
saturation are not established by that check.

\clearpage
\section*{3. Thermal null, physical contradiction, and remaining proof}
The nominal apparent density exceeds the 432 MHz plasma critical density,
$\criticalNe$ m$^{-3}$, by \cutoffFactor{} times. Even the factor 10 lower
stress case is $\lowerNe$ m$^{-3}$, with plasma frequency approximately
\lowerPlasmaMHz{} MHz. Such a value is incompatible with the assumed weak,
propagating, volume-filling thermal ISR interpretation. This is a
\emph{self-consistency rejection of that interpretation}, conditional on the
calibration and gain transfer, rather than a measurement of an overdense
ionospheric layer.

Magnetic perpendicularity alone is not proof of coherent scattering.
Thermal ISR spectra also narrow near perpendicularity, and their collision
model and beam averaging matter \cite{Longley2020}. A performed finite-beam
magnetized-electron BGK sensitivity calculation at $N_e=2.5\times10^{12}$
m$^{-3}$, $T_i=1000$ K and $T_e=1100$ K gives an integrated-power ratio of
\thermalLow--\thermalHigh{} between the event aspect and a $25.4^\circ$
reference. Assumed collision rates span 10--1000 s$^{-1}$ and two-way angular
FWHM 0.5--1$^\circ$. Increasing angular quadrature from 9 to 33 changes the
integrals by at most \angleRefinement\%. The calculation uses the existing
430.2 MHz reference model, versus the actual 432 MHz recording, and a
$\pm50$ kHz ion-line frequency window. Extending the reference spectrum to
the full $\pm2$ MHz recorded band, with refined sampling around electron
gyroharmonics, gives a ratio \fullThermalLow--\fullThermalHigh{}; the
integrated powers change by at most \fullBandChange\%. These are useful
order-of-magnitude thermal
power test, not a fitted event spectrum, a Coulomb-accurate collision model,
or a complete LFM ambiguity calculation. Within these tested assumptions,
thermal magnetic narrowing cannot explain the observed power excess.

The strongest defensible present statement is: \emph{the persistent,
near-perpendicular enhanced echo is strongly inconsistent with a normal
thermal ISR interpretation under the audited calibration, and is consistent
with coherent scattering by field-aligned plasma irregularities.}
At 432 MHz the monostatic Bragg scale is approximately 0.347 m. Coherent
scatter would measure a density Fourier component on that transverse scale;
it would not require the mean density to reach the apparent ISR value.
A stronger mechanism-specific claim should complete three checks:
\begin{enumerate}
\item Fit contemporaneous LP spectra or the full code-aware LFM voltage
covariance with finite-beam magnetized thermal ISR, plus a narrow coherent
component. Predict independent pulses or beams; a narrow spectral peak alone
is insufficient near perpendicularity.
\item Establish the absolute L2 calibration and a contemporaneous independent
density anchor. The local SAO archive checked here is empty; the public GIRO
Sanya SA418 year list has no 2026 entries. A March 2026 ionogram has therefore
\emph{not} been obtained. A valid ionosonde $f_oF_2$ gives
$N_mF_2=1.24044\times10^{10}[f_oF_2/\mathrm{MHz}]^2$ m$^{-3}$ and should be
compared with the calibration profiles, with spread-F scaling uncertainty.
\item Test angular dependence, Doppler, receiver/transmitter coupling and
measured main- and sidelobe patterns. Coded delayed instrumental leakage or
an off-axis scatterer is not ruled out by the wrong-chirp check. Nominal
altitude assumes the echo comes from the intended beam.
\end{enumerate}

\begin{thebibliography}{9}\small
\bibitem{Hao2022} Hao et al. (2022), Initial ionospheric ion line results and
evaluation by SYISR, \emph{JGR Space Physics}, 127, e2022JA030563.
\href{https://doi.org/10.1029/2022JA030563}{doi:10.1029/2022JA030563}.
\bibitem{Yue2023} Yue et al. (2023), On the ion line calibration by plasma line
in ISR measurements, \emph{Remote Sensing}, 15, 1553.
\href{https://doi.org/10.3390/rs15061553}{doi:10.3390/rs15061553}.
\bibitem{Longley2020} Longley et al. (2020), Millstone Hill ISR measurements
of small aspect angle spectra, \emph{JGR Space Physics}, 125, e2019JA027708.
\href{https://doi.org/10.1029/2019JA027708}{doi:10.1029/2019JA027708}.
\end{thebibliography}
\end{document}
